\section{Introduction}\label{sec:intro}

Continuous (implicit) representations decode a signal at arbitrary
coordinates, and a large body of work controls \emph{which frequencies} such a
representation can express: the frequency scale of sinusoidal activations
\citep{sitzmann2020siren,liu2024finer}, the bandwidth of Fourier features
\citep{tancik2020fourier}, target-spectrum-driven hyper-parameter selection
\citep{kania2025fresh,winner2025,sec2026}, per-band capacity allocation
\citep{wieninr}, or explicit spectral priors such as Gaussian processes and
kriging \citep{rasmussen2006gp}. We refer to all of these as \emph{spectral
adaptation}. Its empirical record is uneven. In multi-frame
super-resolution, choosing the Fourier-feature bandwidth well is decisive
\citep{superf2026}; in sparse-pixel reconstruction, a matched spectrum is
worth several decibels (\S\ref{sec:exp-b}); but in single-image
arbitrary-scale super-resolution (ASISR) the same lever barely moves.

ASISR motivates the question. LIIF and LTE decode local features at
continuous coordinates \citep{chen2021liif,lee2022lte}; our existing SC-INR
research implementation uses an analytic footprint readout. Recent
controlled benchmarking reports small gains from more elaborate decoders,
alongside substantial dependence on training recipes
\citep{nasir2026empirical}. Our band measurements show similarity among a
limited collection of trained decoders. This is evidence about those
models and settings, not a proof of architectural saturation.

We use a classical LMMSE reference conditioned on $A$, noise and an assumed
covariance. A generic $1/f^2$ spectrum permits computation before training,
but does not bound all estimators on arbitrary real images. Combined with
exact scalar mismatch costs, it supplies conditional tendencies whose
band-level accuracy and practical selection value must be tested.

\paragraph{Contributions.}
\begin{enumerate}
\item \textbf{A training-free diagnostic.} We cast spectral adaptation as the
choice of a prior spectrum (Prop.~\ref{prop:bochner}) and use the per-band
linear recoverability $\Rlin$ (Prop.~\ref{prop:bound}) to draw an a-priori
three-regime map for single-image, sparse-pixel, multi-frame and 1D signal
tasks (\S\ref{sec:zones}, Table~\ref{tab:map}).
\item \textbf{A closed-form mismatch cost.} Prop.~\ref{prop:mismatch} gives
the excess error of a prior that is off by a factor $\kappa$ in a decoupled scalar channel with
recoverability $R$. Its asymmetry yields a per-band tendency---\emph{wide} is safe as $R\to1$,
\emph{matching} matters at intermediate $R$, \emph{conservative} is safe where
$R$ is low---and its approximate energy-weighted sum over bands empirically ranks GP priors
(though it does not reliably pick the single best); it suggests a qualitative interpretation of noisy grid and Fourier-feature fits, without establishing their mechanism (\S\ref{sec:zones}).
\item \textbf{Pre-registered cross-task validation} (\S\ref{sec:exp}).
The mismatch-cost shape and asymmetry hold on 60 task--band--$\kappa$ cells
(Spearman 0.903 in the independent reimplementation) \src{E2}; the zones move with the ASISR scale factor as
predicted (Spearman 0.952 over 27 cells) \src{E3}; in the low-$R$
regime trained ASISR networks exceed the linear oracle in every cell while
differing among themselves by at most 0.03 in 10 of 12 cells \src{E3},
which describes similarity in these settings; and the benefit of the cell/footprint
readout is negligible on average in the high-$R$ band and in these evaluated seed1 cells appears where the network has
partially recovered the band---consistently in LTE and SC-INR
(cross-architecture Spearman 0.985, seed 1) \src{F1}; averaging the full model
over training scales predicts the ranking of these losses after the fact
(Spearman 0.700--0.954 over four completed seed2/3 groups; correlational) \src{L1L2}.
\item \textbf{A prospective test of practical guidance}
(\S\ref{sec:g2}). Independent-calibration risk selection improves a fixed
GP default on a fully frozen 36-case cohort, but selects the same spectrum
throughout and is less accurate and more expensive than frame-holdout CV.
The test separates a useful calibrated prior from a demonstrated advantage
of target-specific analytical selection \src{G2-result; G2-audit}.
\end{enumerate}

\paragraph{Scope and honesty.} Props.~\ref{prop:bound}--\ref{prop:mismatch}
follow directly from classical Wiener/LMMSE theory; our claim is not new
estimation theory but its use as an a-priori diagnostic and prescription for
spectral adaptation in continuous representations, together with a
pre-registered test across tasks. Several pre-registered criteria failed; we
report them in the main text and in \S\ref{sec:limits}.

\paragraph{Completed replication and scope.} Seed2/3 tests support a simple
footprint predictor in four groups and L2 magnitude in three. Function
similarity fails in all four, and L2 superiority over the simple predictor
fails the 3/4-group criterion. Classical identities, oracle diagnosis and
prospective selection are different claims. G2 supports improvement over
its specified default, but not adaptive superiority or savings over CV.
